\newpage
\section{Scénario 4 : Visualisation d'un (méta)modèle}
\label{sec:scenario4}

\subsection{Description du scénario}

Ce scénario évalue les performances de la boucle de rendu GraSyLa : pour
chaque objet d'un modèle concret, MetaDONE doit déterminer quelle équation
du script s'applique à cet objet, puis évaluer son arbre d'expression pour
produire le rendu graphique.

La pipeline complète de visualisation comprend trois étapes successives :

\begin{enumerate}
  \item \textbf{Récupération des équations} —
        \texttt{GrasylaScript.getEquations()} construit la liste des 106
        équations disponibles.
  \item \textbf{Dispatch} — pour chaque objet du modèle, identifier les
        équations dont le \texttt{MetaObject} du membre gauche (LHS)
        correspond au type de l'objet.
  \item \textbf{Évaluation des expressions} — parcourir et interpréter
        l'arbre d'expression (RHS) de chaque équation applicable pour
        produire les composants graphiques.
\end{enumerate}

Le benchmark implémenté dans \texttt{ScenarioVisualization} couvre les
étapes 1 et 2. L'étape 3 requiert la couche UI Eclipse (SWT, moteur
de rendu SVG) et ne peut pas être isolée dans un test standalone~;
son coût réel est quantifié a posteriori par un profilage live de
l'application Eclipse complète (section~\ref{sec:sc4-live}).

\subsection{Infrastructure de benchmark}

Le script BPMN \texttt{bpmn\_mm.grasyla} (106 équations) est chargé en
mémoire une seule fois en \texttt{@BeforeClass}. Quatre modèles concrets de
taille croissante sont construits programmatiquement :
$N \in \{50, 200, 500, 1\,000\}$ objets répartis en trois types BPMN
(\texttt{Task}, \texttt{ExclusiveGateway}, \texttt{StartEvent}). Ces types
ont été choisis pour leur profondeur de hiérarchie différente, représentative
de l'hétérogénéité d'un vrai modèle.

La mesure centrale est la boucle de dispatch naïve (Listing~\ref{ls-dispatch}).

\begin{lstlisting}[language=Java,
  caption={Boucle de dispatch naïve $O(N \times M)$},
  label={ls-dispatch}]
for (final ConcreteObject co : model.getDefinition()) {
    for (final GrasylaEquation eq : equations) {
        if (co.hasType(eq.getLHS().getMetaObject())) {
            matches++;
        }
    }
}
\end{lstlisting}

Cette boucle de complexité $O(N \times M)$ effectue $N \times M = N \times 106$
appels à \texttt{hasType()}, où chaque appel vérifie si le
\texttt{MetaObject} LHS de l'équation figure parmi les types (direct +
supertypes) de l'objet.

\subsection{Mécanisme de \texttt{hasType()} et rôle du cache JPA}

Lorsqu'un objet est créé via \texttt{createObject(moTask)}, MetaDONE
enregistre dans la table \texttt{NN\_typeof} une entrée pour le type direct
et pour chacun de ses supertypes
(\texttt{Activity}, \texttt{FlowNode}, \texttt{FlowElement},
\texttt{BaseElement}, \texttt{InteractionNode}). Un objet \texttt{Task}
dispose ainsi de $\sim$6 entrées \texttt{NN\_typeof}.

La méthode \texttt{hasType(MetaObject)} dans \texttt{ConcreteObject\_impl}
appelle \texttt{getTarget().getObjectTypes()}, qui itère la collection JPA
\texttt{instanceOf} (\texttt{@OneToMany} sur \texttt{NN\_typeofJPA2}). Cette
collection est chargée depuis la base lors du premier accès puis conservée
dans l'identity map OpenJPA pour toute la durée de la session. Les appels
suivants itèrent un \texttt{HashSet} Java en mémoire, sans aucune requête
SQL.

Il résulte de ce mécanisme que le coût d'un appel à \texttt{hasType()},
bien que théoriquement $O(K)$ avec $K$ le nombre de supertypes, est en
pratique de l'ordre de quelques dizaines de nanosecondes — le coût d'une
itération en mémoire.

\subsection{Résultats de référence}

\begin{table}[h]
\centering
\caption{Résultats du benchmark de dispatch (médiane de 10 répétitions,
mesures complètes en annexe~\ref{ann:mesures-s4}).}
\label{tab:sc4-dispatch}
\begin{tabular}{llrr}
\hline
\textbf{Phase} & \textbf{Paramètre} & \textbf{Médiane} & \textbf{CV} \\
\hline
t1 — \texttt{getEquations()} & 106 équations         & 3,70~ms & 8,4\,\% \\
t2 — Dispatch                & $N = 50$ (650 m.)     & 1,32~ms & 15,8\,\% \\
t3 — Dispatch                & $N = 200$ (2\,600 m.) & 4,00~ms & 8,5\,\% \\
t4 — Dispatch                & $N = 500$ (6\,500 m.) & 8,71~ms & 11,8\,\% \\
t5 — Dispatch                & $N = 1\,000$ (13\,000 m.) & 14,49~ms & 17,3\,\% \\
\hline
\end{tabular}
\end{table}

La constance du ratio $\text{matches}/N = 13$ confirme la cohérence du
benchmark : chaque objet correspond à exactement 13 équations en moyenne,
quelle que soit la taille du modèle.

\subsection{Analyse}

%\paragraph{Croissance sous-linéaire.}
Le rapport des temps entre $N = 50$ et $N = 1\,000$ est de
$14{,}49/1{,}32 \approx 11{,}0\times$, contre $20\times$ attendu pour une
complexité strictement linéaire. Cet écart traduit l'effet du compilateur
JIT de la JVM~: la première itération ($N = 50$, boucle froide) est
exécutée en mode interprété, les suivantes bénéficient d'un code natif
compilé. Le coût par objet passe de 26,4~µs à $N = 50$ à 14,5~µs à
$N = 1\,000$ --- un facteur $\approx 1{,}8$ dû au seul chauffage JIT
(tableau~\ref{tab:sc4-cout-unitaire}).

\begin{table}[h]
\centering
\caption{Coût unitaire du dispatch en fonction de $N$ (calculé à partir
des médianes du tableau~\ref{tab:sc4-dispatch}).}
\label{tab:sc4-cout-unitaire}
\begin{tabular}{rrrr}
\hline
\textbf{N} & \textbf{Temps total} &
\textbf{Coût / objet} & \textbf{Coût / appel \texttt{hasType()}} \\
\hline
50     &  1,32~ms & 26,4~µs & 249~ns \\
200    &  4,00~ms & 20,0~µs & 188~ns \\
500    &  8,71~ms & 17,4~µs & 164~ns \\
1\,000 & 14,49~ms & 14,5~µs & 137~ns \\
\hline
\end{tabular}
\end{table}

%\paragraph{\texttt{getEquations()} plus coûteux que le dispatch sur petits modèles.}
Pour $N = 50$, la récupération des équations (3,70~ms) représente
$\approx 2{,}8\times$ le temps de dispatch (1,32~ms). Cette méthode
reconstruit à chaque appel 106 objets \texttt{GrasylaEquation} en relisant
depuis JPA 4 propriétés et 2 rôles par équation. Si elle est appelée à
chaque cycle de rendu --- ce qui est le cas en pratique --- elle constitue
un surcoût fixe indépendant de la taille du modèle.

%\paragraph{Pertinence des valeurs absolues.}
Un diagramme BPMN réaliste contient 20 à 50 éléments. Pour $N = 50$, le
dispatch complet (\texttt{getEquations} + boucle) prend moins de 5 ms, ce
qui est négligeable dans un cycle de rendu complet. Les optimisations
algorithmiques envisagées (mise en cache de \texttt{getEquations()}, index
\texttt{Map<MetaObject, List<GrasylaEquation>>}) sont techniquement valides
mais ne produiraient pas d'amélioration perceptible sur des modèles de
taille courante.

\subsection{Conclusion du benchmark standalone}

Le benchmark standalone établit que le dispatch d'équations est déjà efficace
dans MetaDONE~: le cache JPA absorbe intégralement le coût des accès type,
ramenant chaque appel à \texttt{hasType()} à $\sim$137~ns à l'état chaud.
Pour des modèles de taille réaliste ($N < 100$), le dispatch total reste
inférieur à 5~ms, ce qui est négligeable dans un cycle de rendu complet.
Le goulot de performance de la visualisation, s'il existe, se situe donc
dans l'évaluation des arbres d'expression (étape~3 du pipeline) et le rendu
graphique --- deux couches inaccessibles en mode standalone et qui font
l'objet de la section suivante.

\subsection{Profilage live de l'application Eclipse}
\label{sec:sc4-live}

Le benchmark standalone (sections précédentes) ne couvre que les étapes 1
et 2 du pipeline de rendu. Pour mesurer l'étape 3 --- l'évaluation
récursive des arbres d'expression et la construction des composants
graphiques --- il faut accéder à la couche UI Eclipse où le moteur GraSyLa
parcourt l'arbre des \texttt{ConcreteObject} liés par
\texttt{bootsMR\_Contains}, résout les références MetaL$_2$, évalue les
functors et compose les composants graphiques. Un profilage par
échantillonnage a donc été réalisé sur l'application Eclipse en cours
d'exécution à l'aide d'\emph{async-profiler}.

\subsubsection{Protocole de mesure}

Un script GraSyLa dédié, \texttt{petri\_benchmark.grasyla}, définit le
métamodèle réseau de Petri (Place, Transition, arcs) et expose une action
\texttt{doubleClick} créant 100 Places et 100 Transitions par appel, via la
méthode \texttt{PetriDemo.createBenchmark()}. Le profilage suit la séquence
suivante :

\begin{enumerate}
  \item lancement d'Eclipse, ouverture d'un modèle et association du script ;
  \item attachement d'async-profiler au processus en cours :
\begin{lstlisting}[language=bash, caption={Attachement d'async-profiler au JVM Eclipse}, label={ls-asprof}]
jps -l | grep eclipse
./asprof start -e cpu -i 5ms <PID>
\end{lstlisting}
  \item cinq doubles-clics successifs (création de $5 \times 200 = 1\,000$
        éléments), entrecoupés de déplacements de quelques n\oe{}uds ;
  \item fermeture des vues, puis arrêt et export du profil :
\begin{lstlisting}[language=bash, caption={Arrêt du profilage et export du flamegraph}, label={ls-asprof-stop}]
./asprof stop -f flame.html <PID>
\end{lstlisting}
\end{enumerate}

Les options \texttt{-XX:+UnlockDiagnosticVMOptions} et
\texttt{-XX:+DebugNonSafepoints} ont été ajoutées à \texttt{eclipse.ini} afin de
permettre la capture de traces de pile précises dans le code compilé par le JIT.
Le profil obtenu totalise \textbf{8\,444 échantillons CPU}.

\subsubsection{Résultats}

Le Tableau~\ref{tab:sc4-live} présente la part de chaque famille de méthodes,
obtenue par recherche par expression régulière dans le flamegraph.

\begin{table}[h]
\centering
\caption{Part du temps CPU par famille de méthodes (8\,444 échantillons)}
\label{tab:sc4-live}
\begin{tabular}{lr}
\hline
\textbf{Famille de méthodes (regexp)} & \textbf{\% CPU} \\
\hline
\texttt{Grasyla} (moteur GraSyLa, toutes méthodes)         & 88,56 \\
\texttt{trigger} (propagation d'événements)                & 53,32 \\
\texttt{buildNewComponent|buildComponent|buildChild}       & 53,21 \\
\texttt{evaluateComponent}                                 & 53,20 \\
\texttt{updateSceneChildren|GraphEngine}                   & 17,91 \\
\texttt{getEquations|hasType} (dispatch, étapes 1--2)      & 11,16 \\
\texttt{precondition\_createRole} (vérification DbC)        & 0,04 \\
\texttt{GroovyClassLoader|parseClass} (compilation script) & 0,05 \\
\hline
\end{tabular}
\end{table}

\textit{Remarque de lecture :} les trois lignes \texttt{trigger},
\texttt{buildComponent} et \texttt{evaluateComponent} affichent toutes
$\approx 53\,\%$ parce qu'elles appartiennent à \emph{la même pile d'appel}
(\texttt{trigger} $\rightarrow$ \texttt{buildComponent} $\rightarrow$
\texttt{evaluateComponent}) ; ces valeurs ne s'additionnent pas.

\subsubsection{Analyse}

%\paragraph{Le moteur GraSyLa domine le profil.}
Avec 88,56~\% du temps CPU, l'évaluation et la construction des composants
graphiques (\texttt{GrasylaEngine}, \texttt{GraphEngine},
\texttt{GrasylaScene}) constituent le goulot d'étranglement réel de la
visualisation. Ce résultat confirme empiriquement l'hypothèse formulée au
Scénario~4 : le coût de rendu réside dans l'étape~3 (évaluation des arbres
d'expression), inaccessible en mode standalone.

%\paragraph{Un rendu piloté par événements.}
La création s'effectuant dans une vue \emph{ouverte}, chaque \texttt{createObject}
émet un événement (\texttt{MetadoneEvents.trigger}) propagé par la tour
d'événements (\texttt{EventTowerNormal.findSubscribers}) jusqu'au gestionnaire
\texttt{GrasylaEngine\$ConcreteObjectChangeHandler}, qui reconstruit la scène.
Ce mécanisme explique les $\approx 53\,\%$ passés dans la pile
\texttt{trigger} $\rightarrow$ \texttt{buildComponent} $\rightarrow$
\texttt{evaluateComponent} : la construction de $N$ objets dans une vue active
engendre une reconstruction incrémentale à chaque insertion. C'est cette boucle
de re-rendu événementiel, et non l'algorithme de dispatch, qui domine le coût
perçu.

%\paragraph{La précondition DbC n'est pas un goulot.}
La vérification \texttt{precondition\_createRole} représente 0,04\,\% du temps
CPU. Ce résultat conforte \emph{a posteriori} le rejet de la
\textbf{candidate \og bypass-précondition \fg{}} étudiée au Scénario~2
(section~\ref{sec:scenario2-precondition-rejet}, voir aussi
chapitre~\ref{chap:discussion})~: le contournement de la précondition aurait
violé le principe \emph{Design by Contract} pour un gain inférieur au bruit
de mesure. De même, la compilation du script Groovy (0,05\,\%) est négligeable,
le script n'étant compilé qu'une fois au chargement.

%\paragraph{Le dispatch reste mesurable en contexte interactif.}
À l'inverse des mesures standalone, où le dispatch était négligeable pour des
modèles de taille réaliste, la famille \texttt{getEquations|hasType} représente
ici 11,16~\% du temps CPU. Cet écart s'explique par le mécanisme de re-rendu :
\texttt{getEquations()} et la sélection d'équation sont réinvoquées à chaque
reconstruction déclenchée par un événement. Ce résultat valide \emph{a
posteriori} la piste d'optimisation à court terme identifiée dans la discussion
(Section~\ref{chap:discussion}) : la mise en cache de
\texttt{GrasylaScript.getEquations()}, invalidée via \texttt{notifyScriptChanged()},
éliminerait un surcoût fixe réinvoqué à chaque cycle de rendu.

\subsection{Conclusion du Scénario~4}

Le Scénario~4 combine deux niveaux d'analyse complémentaires. Le
\textbf{benchmark standalone} a montré que le dispatch d'équations est
déjà efficace en MetaDONE et reste négligeable sur des modèles de taille
réaliste~; il a également exposé sa propre limite, à savoir l'impossibilité
de mesurer l'évaluation des arbres d'expression sans la couche UI Eclipse.
Le \textbf{profilage live} a comblé cette limite et identifié le goulot
réel~: l'évaluation récursive des arbres d'expression GraSyLa concentre
$\approx 88\,\%$ du temps CPU, le coût étant amplifié par le caractère
événementiel du rendu --- toute modification du modèle dans une vue ouverte
déclenche une reconstruction de la scène.

Deux pistes d'optimisation se dégagent, par ordre de rapport gain/effort~:
(i) la mise en cache de \texttt{GrasylaScript.getEquations()}, dont le
surcoût mesuré atteint 11\,\% en contexte interactif, et (ii) une
reconstruction incrémentale ciblée de la scène plutôt qu'une reconstruction
globale à chaque événement. Ces pistes sont reprises dans le
chapitre~\ref{chap:discussion}.

Enfin, la précondition \emph{Design by Contract} de
\texttt{precondition\_createRole} ne contribue pas mesurablement au coût
(0,04\,\%), ce qui valide les arbitrages de conception retenus dans ce
travail et confirme le rejet de la candidate \og bypass-précondition \fg{}.
